% 第5章（正文第12章）图：循环Kaczmarz在贯穿例上的交替投影路径。
% 数据由同目录 data.py 用有理数精确计算：A=[[4,1],[1,3]]，b=(1,2)，x0=0，循环顺序 1,2,1,2,...
% x1=(4/17,1/17)，x2=(67/170,91/170)，x3=(383/2890,679/1445)，x4=(5153/28900,17549/28900)，...；x*=(1/11,7/11)。
% 等比例坐标，使每次投影在图上与对应直线垂直。
\def\DataDir{figures/01-basic-mathematics/03-vectors-matrices-and-multivariate-functions/tikz/c5-kaczmarz-projection-path}
\begin{tikzpicture}
\begin{axis}[width=78mm,axis equal image,
  xmin=-0.04,xmax=0.46,ymin=-0.05,ymax=0.74,
  xtick={0,0.2,0.4},ytick={0,0.2,0.4,0.6},
  xticklabel style={/pgf/number format/fixed},yticklabel style={/pgf/number format/fixed},
  xlabel={$x_1$},ylabel={$x_2$},axis lines=left,clip=false,
  xlabel style={at={(ticklabel cs:1)},anchor=north east,inner sep=1pt},
  ylabel style={at={(ticklabel cs:1)},anchor=south east,rotate=-90,inner sep=1pt}]
  \addplot[FigureMuted,line width=.6pt,domain=0.065:0.26,samples=2] {1-4*x};
  \addplot[FigureMuted,line width=.6pt,dashed,domain=-0.04:0.46,samples=2] {(2-x)/3};
  \addplot[FigureBlue,line width=1.1pt,mark=*,mark size=1.3pt,mark options={fill=white},
    postaction={decorate}] table[x=x1,y=x2] {\DataDir/path.dat};
  \addplot[only marks,mark=x,mark size=2.4pt,FigureInk,line width=.8pt] coordinates {(0.090909,0.636364)};
  \tikzset{lab/.style={font=\FigureNoteFont\fontsize{7.5}{9}\selectfont,inner sep=1pt}}
  \node[lab,anchor=north west] at (axis cs:0.005,-0.005) {$\bm x_0$};
  \node[lab,anchor=west] at (axis cs:0.245,0.055) {$\bm x_1$};
  \node[lab,anchor=north west] at (axis cs:0.40,0.525) {$\bm x_2$};
  \node[lab,anchor=east] at (axis cs:0.125,0.47) {$\bm x_3$};
  \node[lab,anchor=south west] at (axis cs:0.182,0.612) {$\bm x_4$};
  \node[lab,anchor=south east] at (axis cs:0.085,0.65) {$\bm x_\star$};
  \node[lab,text=FigureMuted,anchor=south] at (axis cs:0.065,0.745) {方程1};
  \node[lab,text=FigureMuted,anchor=south east] at (axis cs:0.46,0.555) {方程2};
\end{axis}
\end{tikzpicture}
